\documentclass[a4paper,11pt]{article}
\usepackage[utf8]{inputenc}
\usepackage{hyperref}
\usepackage{listings}
\usepackage{color}
\usepackage{geometry}
\usepackage{tcolorbox}

\geometry{a4paper,margin=2.5cm}

\definecolor{codegreen}{rgb}{0,0.6,0}
\definecolor{codegray}{rgb}{0.5,0.5,0.5}
\definecolor{backcolour}{rgb}{0.95,0.95,0.92}

\lstdefinestyle{mystyle}{
    backgroundcolor=\color{backcolour},   
    commentstyle=\color{codegreen},
    keywordstyle=\color{blue},
    numberstyle=\tiny\color{codegray},
    stringstyle=\color{codegreen},
    basicstyle=\ttfamily\footnotesize,
    breakatwhitespace=false,         
    breaklines=true,                 
    captionpos=b,                    
    keepspaces=true,                 
    numbers=left,                    
    numbersep=5pt,                  
    showspaces=false,                
    showstringspaces=false,
    showtabs=false,                  
    tabsize=2
}

\lstset{style=mystyle}

\begin{document}

\title{COSMIC Installation Guide}
\author{COSMIC Development Team}
\date{\today}

\maketitle

\section{Introduction}
COSMIC (Compact Object Synthesis and Monte Carlo Investigation Code) is a rapid binary population synthesis suite with a special focus on generating compact binary populations. This guide will walk you through the installation process.

\section{Prerequisites}
Before installing COSMIC, you need to have certain prerequisites installed on your system:

\subsection{Installing gfortran}
Since COSMIC requires compilation of Fortran code, you'll need a gfortran installation. The installation method depends on your operating system:

\subsubsection{For MacOS Users}
You can install gfortran using Homebrew:
\begin{lstlisting}[language=bash]
brew install gcc
\end{lstlisting}

\begin{tcolorbox}[title=Note for Apple Silicon Users]
If you're using an Apple Silicon (ARM) processor, verify your Python architecture:
\begin{lstlisting}[language=bash]
python -c "import platform; print(platform.architecture())"
\end{lstlisting}
If the output shows \texttt{('64bit', 'x86\_64')}, install the ARM version:
\begin{lstlisting}[language=bash]
brew install python@3.10
\end{lstlisting}
\end{tcolorbox}

\subsubsection{For Unix/Linux Users}
Install gfortran using your package manager:
\begin{lstlisting}[language=bash]
sudo apt-get install gfortran
\end{lstlisting}

\subsubsection{For Windows Users}
Windows is not directly supported due to issues with the gfortran compiler and libraries. However, you can use Windows Subsystem for Linux (WSL):
\begin{enumerate}
    \item Install WSL by following Microsoft's guide at: \url{https://docs.microsoft.com/en-us/windows/wsl/install}
    \item Follow the Unix/Linux instructions within WSL
\end{enumerate}

\section{Installing COSMIC}
We recommend using conda to create a dedicated environment for COSMIC. This ensures all dependencies are installed correctly.

\subsection{Setting up the Conda Environment}
First, install Miniconda if you haven't already from: \url{https://docs.conda.io/en/latest/miniconda.html}

\subsubsection{For MacOS and Unix/Linux}
\begin{lstlisting}[language=bash]
conda create -n cosmic numpy h5py python=3.10
source activate cosmic
pip install cosmic-popsynth
\end{lstlisting}

\section{Using Jupyter with COSMIC}
To use COSMIC with Jupyter notebooks, you need to install Jupyter in the COSMIC environment:

\subsection{Using conda}
\begin{lstlisting}[language=bash]
conda install jupyter ipython
\end{lstlisting}

\subsection{Using pip}
\begin{lstlisting}[language=bash]
pip install jupyter ipython
\end{lstlisting}

\begin{tcolorbox}[title=Important Note]
Always use Jupyter and IPython from within the COSMIC conda environment. Using the global instance may cause issues.
\end{tcolorbox}

\section{Troubleshooting}
\begin{itemize}
    \item \textbf{MacOS Users}: If you encounter issues, try reinstalling gfortran and then reinstall COSMIC
    \item \textbf{Linux Users}: Make sure your gfortran installation is up to date
    \item \textbf{WSL Users}: Ensure you're using WSL 2 and have installed all required Linux components
\end{itemize}

\section{Next Steps}
After installation, you can verify your installation by running a simple test:
\begin{lstlisting}[language=python]
from cosmic.sample.initialbinarytable import InitialBinaryTable
from cosmic.evolve import Evolve
\end{lstlisting}

If no errors appear, your installation is successful!

\end{document}